\section{The full split-prime criterion for the three quadratic cubic twists}
\label{sec:mixed-cubic-split-cube-quadratic}

Put $K=\mathbb Q(\omega)$, $t^3=269$, $L=K(t)$, and let
$\chi$ be either faithful character of $\operatorname{Gal}(L/K)$.
For $d\in\{-1,2,-2\}$ set
\[
 \begin{gathered}
 L_d=L(\sqrt d),\qquad B_d=\mathbb Q(t,\sqrt d),\\
 (\delta_{-1},\delta_2,\delta_{-2})=(\chi_{-4},\chi_8,\chi_{-8}).
 \end{gathered}
 \tag{MCN.1}
\]
The following expands the split-prime scope of AC-G207 using
its already certified class inputs.  No new class computation
or finite character is introduced.

For each $d$, define
\[
 \begin{split}
 \mathscr P_d^* =\{\ell\text{ prime}:{}&\ell>3,\quad\ell\ne7,269,
                    \quad\ell\equiv1\pmod3,\\
 &269^{(\ell-1)/3}\not\equiv1\pmod\ell
                  \ \text{or}\ \delta_d(\ell)=-1\}.
 \end{split}
 \tag{MCN.2}
\]
For $\ell\in\mathscr P_d^*$ let $O$ be the original ring of
integers of any finite extension $E/\mathbb Q_\ell$ containing
$\mu_3$, with uniformizer $\varpi$ and residue field $k$.
Let $\kappa:G_K\to1+\mathfrak m_O$ be continuous and
unramified away $\ell$.  For $j=1,2$ and a positive integer
$e$ divisible by $\ell-1$, put
\[
             D=(E/O)(\chi^j\delta_d\kappa^{-1}c_\ell^{-e}).
 \tag{MCN.3}
\]
Full strict means zero restriction to $H^1(K_v,D)$ at both
$v\mid\ell$ and to $H^1(I_v,D)$ at every other finite place;
use these same conditions for $D[\varpi]$.

\begin{theorem}[AC-G215: full split-prime quadratic criterion]
\label{AC-G215}
For all the choices in (MCN.1)--(MCN.3),
\[
       \mathcal S_\ell(K,D[\varpi])=0,
       \qquad \mathcal S_\ell(K,D)=0.
 \tag{MCN.4}
\]
The part added to AC-G207 is described by the prime set
\[
 \begin{split}
 \mathscr C=\{\ell\text{ prime}:{}&\ell>3,\quad\ell\ne7,269,
                          \quad\ell\equiv1\pmod3,\\
 &269^{(\ell-1)/3}\equiv1\pmod\ell,
                          \quad\ell\not\equiv1\pmod{24}\}.
 \end{split}
 \tag{MCN.5}
\]
For every integer $T\geq2$, define its actual new triples by
\[
 \begin{split}
 \mathcal C(T)=\{(\ell,\chi^j\delta_d,e):{}&\ell\in\mathscr C,
             \ j=1,2,\ d\in\{-1,2,-2\},\\
 &\delta_d(\ell)=-1,\quad1\leq e<T,\quad\ell-1\mid e\}.
 \end{split}
 \tag{MCN.6}
\]
All labels occur in the original host (FG.3), with their
rational conductor-seven and auxiliary $q,d$ factors trivial.
The auxiliary source parameter $d$ is distinct from the fixed
squarefree integer in (MCN.1).
They retain exactly the original coefficient order, rank
$b_\ell$, and full local conditions (FC.1)--(FC.2).
There are exactly
\[
       |\mathcal C(T)|=4\sum_{\substack{\ell\in\mathscr C\\\ell\leq T}}
                         \left\lfloor\frac{T-1}{\ell-1}\right\rfloor
 \tag{MCN.7}
\]
such triples.  They are disjoint from the noncube-prime
family of AC-G207 and the coefficient-seven family of AC-G211.
For every retained set $J$ of actual triples in (CP.3),
\[
                  \mathcal R_J=\mathcal R_{J\setminus\mathcal C(T)}.
 \tag{MCN.8}
\]
In particular $37\in\mathscr C$ and its four new labels are
$\chi\chi_8$, $\chi^2\chi_8$, $\chi\chi_{-8}$, and
$\chi^2\chi_{-8}$, at all positive weights divisible by $36$.
This entire family is proved at once; the displayed example
does not replace the prime-set assertion.
\end{theorem}

\begin{proof}
\proofstep{1. The inherited class inputs apply before any local cube test.}
The complete AC-G207 field and certificate proof gives
\[
 \begin{array}{c|ccc}
 d&-1&2&-2\\\hline
 \operatorname{Cl}(B_d)\text{ is a quotient of}&C_7&C_{63}&C_{42}.
 \end{array}
 \tag{MCN.9}
\]
These are unconditional completed flag-one certificates,
with independent executions retained there.  They require
neither a new execution nor equality of actual class numbers.
For every prime $\ell>3$, $\ell\ne7$, (MCN.9) gives
$\operatorname{Cl}(B_d)\{\ell\}=0$ independently of the
value of $269$ modulo $\ell$.

The same fields satisfy $\operatorname{Gal}(L_d/K)=C_3\times C_2$.
Let $a$ generate the cubic factor and let $\tau$ fix $t,\sqrt d$
and invert $\omega$.  Then $L_d^{\langle\tau\rangle}=B_d$
and $\tau a\tau^{-1}=a^{-1}$.
On $C=\operatorname{Cl}(L_d)\{\ell\}$, the exact identities
$N\circ j=2$ and $j\circ N=1+\tau$ identify
$C^{\tau=1}$ with $\operatorname{Cl}(B_d)\{\ell\}=0$.
Thus $\tau=-1$ on $C$.  Consequently $a=a^{-1}$ on $C$;
as $a^3=1$, $a$ acts trivially there.
This kills the required nontrivial cubic class components,
without asserting that the whole group $C$ is zero.

The residual coefficient is
$D[\varpi]=k(\overline\chi^{\,j}\overline\delta_d)$:
$\kappa$ and $c_\ell^e$ both reduce to one, and reduction
is injective on $\mu_6$.  This residual character has order
six, and its splitting cover $L_d/K$ has degree six.
Restriction on global $H^1$ is injective by averaging over
this prime-to-$\ell$ group.  A full strict residual class
restricts to an unramified additive character of $G_{L_d}$
and hence, by class field theory, to an equivariant map
$f:C\to k$.  Equivariance gives
$f(ax)=\overline\chi^{\,j}(a)f(x)$.
Since $a=1$ on $C$ and its displayed scalar is nontrivial,
$f=0$.  This proves the first assertion of (MCN.4).

\proofstep{2. Both split coefficient places, including the cube case.}
The prime $\ell$ splits in $K$, so both completions are
$\mathbb Q_\ell$.  Both $\chi^j$ and $\delta_d$ are
unramified there because $\ell\nmid2\cdot3\cdot269$.
If $269$ is a noncube modulo $\ell$, the unramified cubic
Frobenius at either place has order three.  Multiplication
by the quadratic sign $\delta_d(\ell)$ cannot make that
value one in characteristic $\ell>3$.

If $269$ is a cube modulo $\ell$, the three roots of
$X^3-269$ are distinct modulo $\ell$, since $\mu_3\subset
\mathbb F_\ell$ and $\ell\nmid3\cdot269$.
Hensel lifting puts all three roots in $\mathbb Q_\ell$.
Thus $\chi$ is trivial on each of the two local groups.
Condition (MCN.2) in this case says exactly
$\delta_d(\ell)=-1$.  Hence at both places
\[
       \overline\chi^{\,j}(\operatorname{Frob}_v)
       \overline\delta_d(\operatorname{Frob}_v)=-1\ne1.
 \tag{MCN.10}
\]
The residual scalar of any Frobenius lift is therefore
nontrivial in either case, even if $\kappa$ is ramified.
Its scalar minus one is a unit of the original $O$,
so $D^{G_{K_v}}=0$ for both $v\mid\ell$.

At every $v\nmid\ell$, $\kappa$ and $c_\ell$ are unramified.
The remaining inertia image has order dividing six, prime
to $\ell$.  Its invariants on $D$ are zero if this image
is nontrivial and are the divisible module $D$ otherwise.
Every prescribed local invariants quotient modulo $\varpi$
is therefore zero, including the places over $2,3,269$.

\proofstep{3. Original-uniformizer lifting.}
For a global or prescribed local group $G$, the coefficient
sequence gives the exact part
\[
 D^G/\varpi D^G\longrightarrow H^1(G,D[\varpi])
       \longrightarrow H^1(G,D)[\varpi]\longrightarrow0.
 \tag{MCN.11}
\]
A global strict class killed by $\varpi$ has a residual
lift.  Each local restriction of the lift maps to zero,
and Step 2 makes the first term of (MCN.11) zero locally.
The lift is consequently full strict and vanishes by
Step 1.  The original strict group has no $\varpi$-torsion.
Every continuous cocycle into the discrete $E/O$ has finite
image, so its class is killed by some power of $\varpi$.
Any nonzero strict class would yield nonzero
$\varpi$-torsion.  This proves the second assertion of
(MCN.4) on the original ring, with all local conditions.

\proofstep{4. The complete new family and its exact count.}
For odd primes different from three, the split condition
leaves the residue classes $1,7,13,19$ modulo $24$.
The elementary quadratic character formulas give
\[
 \begin{array}{c|rrr}
 \ell\bmod24&\chi_{-4}(\ell)&\chi_8(\ell)&\chi_{-8}(\ell)\\\hline
 1&1&1&1\\
 7&-1&1&-1\\
 13&1&-1&-1\\
 19&-1&-1&1
 \end{array}
 \tag{MCN.12}
\]
Thus a split cube prime supplies no new label in residue
class one and exactly two negative quadratic signs in
each other class.  Both nontrivial cubic characters can
be used for each sign, giving exactly four labels at
every prime of (MCN.5).  The source host has $h\geq1$ and
contains $\mathbb Q(\zeta_8)$; the three quadratic
characters remain distinct over $K$, and the cubic
factor is independent by (FC.1).

Their values already lie in the original completion
$O_\ell$ of $\mathbb Z[\zeta_{3\cdot2^h}]$.
Applying (MCN.4) with $\kappa=1$ and inflating the actual
$G_{K,\Sigma_\ell}$ strict classes into $G_K$ preserves
all full local conditions and the same normalizing $b_\ell$.
There are exactly $\lfloor(T-1)/(\ell-1)\rfloor$ positive
eligible weights, zero unless $\ell\leq T$; this proves
(MCN.7).  The cubic factor is nontrivial, so these are
the full original mixed summands of (CP.3).  Removing
their zeros from any retained $J$ proves (MCN.8) without
an additional comparison cost or a new choice of window.

Finally $269\equiv10\pmod{37}$ and $10^{12}\equiv1\pmod{37}$.
Since $37\equiv13\pmod{24}$, the two negative signs are
exactly $\chi_8$ and $\chi_{-8}$.  This verifies all four
displayed original labels and their $36$-divisible weights.
\end{proof}

\paragraph{Finite-cover dictionary and remaining scope.}
The finite original coefficients $D[\varpi^n]$ are locally
constant etale sheaves on
$\operatorname{Spec}\mathcal O_K[1/(2\cdot3\cdot269\cdot\ell)]$.
The same $C_3\times C_2$ cover $L_d/K$ trivializes the
residual action; the strict completion and inertia kernels
retain the original coefficient modules.  This is the
retained SGA~1, Expose V, sections 4 and 7 dictionary,
with no fresh archival or motivic realization claim.
The author proposed the split-cube/negative-sign extension;
the root independently checked the route and requested
its complete prime-family statement.  The G207 class
certificates and their original execution ownership are
unchanged.  No strict conclusion for inert coefficient
primes, no individual derived-index zero, and no closure
of the all-prime remainder or H04.4 is asserted.
